% \begin{savequote}[75mm]
% Nulla facilisi. In vel sem. Morbi id urna in diam dignissim feugiat. Proin molestie tortor eu velit. Aliquam erat volutpat. Nullam ultrices, diam tempus vulputate egestas, eros pede varius leo.
% \qauthor{Quoteauthor Lastname}
% \end{savequote}

\chapter{General Discussion}

\section{Experimental Considerations for Studying JONs}

In this study, we developed an \textit{in situ} preparation that allows us to target \textit{Drosophila} JONs with patch-clamp electrodes. This experimental approach has strengths and weaknesses that complements other available methods. There are tradeoffs between cellular specificity, resolution (both in time and signal), and throughput. We first consider other electrophysiological techniques, then imaging techniques, and finally JON patch-clamp in a laser dissected organ.

\subsection{Comparison to Electrophysiological Alternatives}

In general, electrophysiological techniques have superior time resolution, and this is true of the most widely used population activity measure for studying JONs: antennal field potential recordings (AFPs). AFPs are sometimes more operationally called sound evoked potentials, or, in reference to vertebrate auditory nerve fiber recordings, compound action potentials \citep{Eberl2000-db,Eberl2011-qa}. The technique is relatively straightforward, and requires only inserting an electrode into a small gap between the antenna and head capsule. This leaves the organ completely intact, and the a3/a2 joint free to rotate in response to mechanical stimulation. 

The AFP is made up of the combined spiking and subthreshold signals from many JONs, and is typically idealized as a full population recording \citep{Eberl2011-qa}. However, the position of a high impedance electrode within the nerve will bias the signal to different subsets of JONs, which project in a topographic manner from JO into the brain \citep{Kamikouchi2006-nj}. AFPs can also be measured by impaling the nerve from within the dissected head capsule. This is more involved, as the thick sheath around the antennal nerve can be difficult to penetrate. With care however, this allows more precise placement of a high impedance electrode next to fluorophore-tagged axons. Those axons could also be tagged with an excitatory or inhibitory channelrhodopsin to confirm class identity \citep{Govorunova2015-fh,Klapoetke2014-bx}, as is often done in mammalian extracellular recordings \citep{Buzsaki2015-yy}. Still, signals recorded in this manner will represent the summed membrane voltage activity of at least dozens of JONs. AFP recordings have yielded several interesting observations, including the basic phenomena of frequency doubling, mean adaptation, and variance adaptation \citep{Clemens2018-wp,Lehnert2013-cl,Tootoonian2012-qc,Jezzini2018-gz,Pezier2013-ef,Clemens2021-mx}.

Better resolution can be obtained by making voltage clamp recordings from a neuron that is gap-junctionally coupled to JONs, namely the giant fiber neuron (GFN) \citep{Lehnert2013-cl}. By adding TTX to the bath, this further permits recording the subthreshold mechanically-evoked generator currents (i.e., pre-spiking transduction currents) from the JONs that are electrically coupled to the GFN. These are primarily A and B-JONs \citep{Matsuo2016-sx,Pezier2014-ho,Lehnert2013-cl}. GFN recordings have been used to demonstrate that elimination of the TRPV channels Nan or Iav abolishes all measurable generator current in JONs, whereas elimination of the TRPN channel NompC decreases the slope of the displacement-current relationship \citep{Lehnert2013-cl}.

To make either AFP or GFN recordings more selective, transduction can be limited to a subset of JONs using a rescue technique. In a TRPV null mutant background (\textit{iav\textsuperscript{1}}) with wild-type \textit{iav} expression in a subset of JONs specified by a Gal4 line \citep{Chang2016-lf,Lehnert2013-cl,Kwon2010-nv}. Practically, however, AFP and GFN measurements still are not able to resolve signals from single cells.
 
There are other electrophysiological techniques that could in principle permit access to smaller subsets of \textit{Drosophila} JONs. For example, suction electrode recording is a straightforward and typically stable recording technique for measuring extracellular multi-unit activity in insect nerves \citep{Miller1979-ye}. Because JON axons enter the brain in large bundles within the antennal nerve, suction electrode recordings are possible. In this technique one draws up the bouquet of severed axonal endings into a relatively large-bore pipette. We found that this method yields higher SNR than AFP recordings if it is targeted to fluorophore-tagged subregions of the nerve (data not shown). JON subtype identity then requires tagging with an excitatory or inhibitory channelrhodopsin as before \citep{Govorunova2015-fh,Klapoetke2014-bx}. In fact, for small pipettes (inner diameter ~0.1-0.2 $\mu m$), we found it was sometimes possible to isolate single unit activity in such a configuration (data not shown). However, these recordings were difficult to interpret, due to the small SNR and high similarity among spike waveforms from different axons. For example, any recording which passes typical extracellular single-unit quality checks \citep{Quiroga2004-rz,Hill2011-zl} is still very likely to represent the combined activity of several, sparsely tuned neurons.

Most published single-cell electrophysiological recordings from chordotonal neurons were performed with sharp glass electrodes \citep{Field1998-hv}. Sharp electrodes have been used to record from JON axons in the antennal nerve axons of locusts \citep{Uchiyama1956-ep} and more recently in hawkmoths \citep{Dieudonne2014-ru}. In blowflies, it is possible to obtain sharp electrode recordings from JON somata or dendrites by inserting the electrode blindly into Johnston’s organ \citep{Burkhardt1960-kk,Schlegel1970-ry}.

Of course, all of these animals are much larger than \textit{Drosophila}. For example, a pair of \textit{Drosophila} antennae would fit side-by-side within the hawkmoth antennal nerve, and an entire \textit{Drosophila} would fit inside a pair of \textit{Calliphora} antennae. Nonetheless, it might be feasible to extend sharp glass microelectrode recordings to  \textit{Drosophila} JON axons, given that some of these axons are relatively large, approximately 1-2$\mu m$ in diameter \citep{Sivan-Loukianova2005-ew}. Indeed, after some optimization, we found it was possible to occasionally record stable spiking responses with sharp electrodes from the axons of the antennal nerve (Supplemental Figure \ref{fig:fig_01s}). However, these recordings were quite brief by comparison to those in larger insects. This is perhaps not surprising given the relatively small diameter of the \textit{Drosophila} antennal nerve when compared to larger insects. Moreover, we could not fill neurons for post hoc identification, which may also be expected given that it is often difficult to fill from the soma even with whole-cell recordings in the central brain. Although there are some configurations that may improve the throughput of this technique, it is difficult to imagine widespread use among \textit{Drosophila} JON researchers.

\subsection{Comparison to Functional Imaging}

By contrast to these electrophysiological techniques, functional imaging has been widely adopted in the study of \textit{Drosophila} JONs. Calcium imaging in particular, using genetically encoded calcium indicators (GECIs), has proved extremely useful in this regard. This technique was first used in a pair of landmark papers \citep{Kamikouchi2009-mz,Yorozu2009-hc}, which showed distinct tuning in different anatomical types of JONs. Specifically, the JONs that project to the A and B branches of the antennal nerve were found to respond preferentially to vibrations; meanwhile, the JONs that project to the C and E branches responded preferentially to static displacements. A followup paper observed that the JONs projecting to the D branch have mixed selectivity for antennal vibration and static deflections \citep{Matsuo2014-up}. Yet another class of JONs specifically implicated in grooming responses was discovered shortly afterwards \citep{Hampel2015-iq,Hampel2020-uw}. Finally, calcium imaging was used to demonstrate that the loose tonotopic map in JON axons is preserved and reorganized at the level of downstream neurons \citep{Patella2018-tk}. 

Some of these calcium studies imaged JON axonal projections into the brain \citep{Yorozu2009-hc,Matsuo2014-up,Patella2018-tk,Hampel2015-iq,Hampel2020-uw}, whereas others imaged JON somata in the antenna \citep{Kamikouchi2009-mz}. Imaging along the antennal nerve is also possible, but signal separation is challenging. Imaging in the brain is by far the most common location, it has the advantage that the antenna may be free for natural stimulation, for example by wind \citep{Patella2018-tk}. Moreover the animal can perform tethered behavior during brain imaging, for example JON activity can be monitored during tethered flight \citep{Mamiya2015-wg}. This represents a large advantage over electrophysiological techniques where movement is problematic. Further, imaging in the brain with selective driver lines is sometimes capable of near single-neuron resolution in the axons \citep{Ishikawa2017-ew}. This technique has also aided understanding of the femoral chordotonal organ neurons in the legs, which have axons in the ventral nerve cord that can be imaged \citep{Mamiya2018-iv}.

The alternative recording location is directly in the organ, where JONs have strong calcium responses \citep{Kamikouchi2009-mz,Kamikouchi2010-jm}. This has been less attractive as the somata are far closer together than their differential projections into the brain, and resolving individual somata during functional imaging in the antenna can be challenging. The thick cuticle, and nonuniform, striated internal structure of the organ make a2 a strongly scattering imaging location. For these reasons driver lines have been required to isolate subtype-specific responses, and organ imaging studies use averaged calcium activity like in the axonal imaging case. We found that scattering can be significantly reduced by cuticle laser dissection, resulting in both increased resolution and SNR (data not shown). Nevertheless, individual somata are large enough to be resolved in the intact organ with standard imaging wavelengths (as evidenced by confocal images), and specialized imaging systems \citep{Katona2012-rr,May2021-ki} or ever increasing sensor SNR \citep{Zhang2020-yz}, may circumvent these problems.

In addition to GECIs, genetically encoded voltage indicators (GEVIs) have recently been used successfully in flies \citep{Yang2016-ef,Chamberland2017-wg,Wienecke2018-ij}, although not yet JONs (but see \citep{Kay2016-xv}). Some GEVIs permit single action potential waveform resolution (“fast” GEVIs), whereas others typically cannot (“slow” GEVIs). To date, all GEVIs used to systematically characterize neural responses in \textit{Drosophila} have relied on slow GEVIs. The relationship between calcium signal in the axons and the firing rate responses at the soma is not known, and it is likely that measuring population activity of JONs in the axons with GEVIs may yield a finer resolution understanding of response diversity. For the reasons discussed above, imaging the relatively low SNR GEVIs in somata would be challenging.

In the future, it is possible that fast GEVIs could be useful to study JONs. However, they have even lower SNR and currently require a high level of expression. This becomes problematic as adding more voltage sensitive domains along the membrane adds a significant capacitance—so called capacitive loading \citep{Yang2016-em}. This phenomenon can be devastating for a neuron's current-to-firing-rate relationship, particularly in small neurons which will require a high density of expression to detect voltage-dependent fluorescence changes. Conventional wisdom holds that the slower GEVIs may be expressed at lower levels which may avoid this issue \citep{Yang2016-em}. 

However, even with this caveat, GEVIs may still prove useful. For example, comparisons could be made between neurons each expressing similar levels of a given indicator. Or, if imaging inside Johnston’s organ, experiments could focus solely on transduction by blocking spikes. This avoids a major problem of identifying spikes in often low SNR recordings and permits trial-averaging the transduction-driven voltage changes. Moreover this may be one of the only ways to investigate the spatial nature of transduction in \textit{Drosophila}, where minimally disruptive sharp electrode access to distinct regions (as done in larger insects \citep{Oldfield1986-xu}) would be difficult.

\subsection{\textit{In situ} Patch-Clamp}

Repeatable, targeted access to individual neurons with patch-clamp has been noted before as an essential step forward towards a mechanistic understanding of JON responses \citep{Eberl2016-yt}. We are not the first investigators to attempt these recordings in \textit{Drosophila}. Clemens and colleagues manually dissected a2 and made recordings from single JONs, but found they were mechanically insensitive \citep{Clemens2018-wp}. Some investigators have also had success with patch clamp recordings in \textit{Drosophila} larval chordotonal organ neurons \citep{Zhang2013-gm}. Notably, a recent study successfully used patch-clamp electrodes to record stimulus-evoked activity in individual neurons of an auditory organ in the locust using an \textit{in situ} preparation \citep{Warren2018-ls}. 

By using patch-clamp in the laser dissected organ, we were able to make recordings from enough JONs to see a variety of responses to antennal stimulation. The number and diversity of recorded neurons is similar to those datasets collected by earlier sharp electrode studies in the femoral chordotonal organ neurons, e.g. in stick insects \citep{Hofmann1985-ny}. One notable advance, however, is that patch-clamp recording can be targeted, with fluorescent tags, to anatomically defined subsets of neurons. We selected A- and B-JONs in part because they represent the largest subset, making targeting easier. However, with alterations to the angle and depth of laser dissection, the preparation is compatible with recording from any subset of JONs. Another reason we selected A- and B-JONs is they were known to be spiking \citep{Lehnert2013-cl}, which was required for cell-attached recordings. 

Whole-cell recordings are a useful complement to cell-attached recordings, because they provide insight into the voltage fluctuations that drive spiking. Whole-cell recordings from chordotonal neurons in a locust auditory organ were recently reported \citep{Warren2018-ls}. We were able to obtain a small number of whole-cell current-clamp recordings, although these were much more difficult to obtain than cell-attached recordings, and so they do not comprise a large portion of our data set. The similarity of JON responses in cell-attached and whole-cell mode is evidence that whole-cell recordings do not grossly disrupt mechanosensitivity.

In sum, there are a variety of methods which are useful for studying \textit{Drosophila} JONs. AFP and GFN recordings have high time resolution, and give the possibility of modest subtype specific targeting. More niche techniques like suction electrode recording may further refine that access. GECI imaging has been a dominant, and useful technique. It is compatible with both population and single neuron recording, however the latter has not yet been shown on a large scale. With some methodological work, GEVIs may provide further, more targeted insights. And finally, the method we have developed, which although somewhat involved, grants targeted single cell access. This diversity of approaches will undoubtedly help advance our understanding of JON physiology, becoming more powerful as genetic tools to target and manipulate JONs are being continually improved \citep{Hampel2020-uw}.

\section{Towards a Computational Framework}

A major component of our study has been an effort to construct a modeling framework for understanding the various types and features of JON responses. Our approach has been to suggest a series of stages which divide up the transformation from stimulus to response, as an attempt to create an interpretable, generative model. A cascade of computational elements is common in many models of peripheral processing, including for JO \citep{Clemens2018-wp,Clemens2021-mx}, although incorporating biophysical features into the stages is somewhat less common. Cascades with biophysical components have been used in the mammalian auditory periphery \citep{Sumner2002-fq}, and somatosensory system \citep{Lesniak2014-kt,Bell1992-yn}, and retina \citep{Schroder2019-ps}. As in other models there were several simplifications required, in our case to bridge antennal motion and JON membrane voltage. We consider each component a hypothesis generating tool, which can help to inform future experiments. Future experiments will likely also prompt a revision of these models. Below we consider each stage, its utility, limitations, and possible extensions.

\subsection{Geometrical Organ Model}

At the first stage of our cascade, we transform the rotational motion of the antenna about the a2/a3 joint into displacements of individual scolopidia. The goal of this stage is to simplify the understanding of JON responses by formalizing a relationship between this rotation and stretch or compression in individual scolopidia. There are a few advantages of this approach. First, it mathematically separates the “input nonlinearity” from subsequent computations in cascade models \citep{Ahrens2008-bb,French1989-rz}. This provides a convenient conceptual division between conversion of the stimulus space (the input nonlinearity) and subsequent stages that deal more with integrating similarity of that stimulus in time. In our model, this stage essentially creates a new frame of reference for mechanoelectrical transduction. A second advantage is that defining this relationship forces us to be explicit about which geometrical components of the antenna are important. Due to the absence of a detailed biophysical model of Johnston's organ, we cannot tie parameters to concrete understanding of force transmission. Instead, we construct the relationship through phenomenological and anatomical observations.

The prevailing geometrical model of \textit{Drosophila} JO is that the point of attachment, the “hook”, curves into the JON array rather than extending straight up along a rotational axis. This has been compared to the action of the bit of a key in a lock \citep{Gopfert2001-fu,Gopfert2002-cw}. In this configuration, scolopidia which are most anterior, or most medial (as defined by \citep{Kamikouchi2006-nj}) attach nearly perpendicularly to the bit of this key, and are therefore displaced nearly linearly with rotation. Those which attach closer to parallel the bit have a more quadratic relationship between the antennal movement and scolopidial length. The exact nature of these attachments and the true axis of rotation, however, are unclear \citep{Pezier2014-ho}. 

We assume that all JONs are attached along the plane of this rotating element. However, this simplifying assumption is only strictly valid for those scolopidia attached radially about the axis of rotation, like spokes on a wheel. JO is instead actually more of an inverted “bowl-shaped” organ, and so to account for this, a second (elevation) angle is required. For slight deviations from the plane, this discrepancy results in only minor asymmetries for motion in either direction, as in the simplified version, but more extreme elevation angles are likely however.

What other parameters might be required of a full biophysical model? Some authors have suggested JONs are only alternately stretched with antennal motion in either direction \citep{Gopfert2002-cw,Albert2007-km}, meaning that no JONs are stretched bidirectionally. For this to be true at any attachment angle in our model, it would require adding an element which prevents compression, similar to the “slackening element” or “catch” in hair cell models \citep{Nadrowski2009-yy}. Additionally, for some attachment angles, it might be possible for the scolopidia to begin to “bend” around the rotational element. In this situation, scolopidia length would be governed partially by an arclength around the element, rather than the simple trigonometric relationship we derived. Finally, a full model would incorporate viscoelastic effects prior to the scolopidium; we have ignored these effects, allowing us to reduce this stage to a point nonlinearity.

How might we obtain data to constrain a fuller biomechanical model? Fortunately, it is possible to fluorescently tag different components of the scolopidia and record their motion during manipulation \citep{Prahlad2017-tx}. Prahlad and colleagues did exactly this for the larval chordotonal organ neurons \citep{Prahlad2017-tx}. Attachment cells, cap cells, and JONs may all be independently tagged. Using 2-photon volumetric imaging, the entire organ can be imaged. Moreover, to improve signal within the organ, cuticle can be thinned using a DUV laser dissection system. In the simplest version of this experiment, we would measure the steady-state position of all components in the intact organ for different antennal positions. A more sophisticated approach using a higher-speed, or more precisely timed, imaging technique could potentially capture viscoelastic effects. These experiments would allow us to build a more biophysically accurate model by constraining the values for attachment angle, and length changes, for example. Then, by acquiring high-resolution images of neurons recorded during electrophysiology experiments, the strong predictions made by the attachment could be tested (e.g. the phases of activation). Finding more accurate parameter values at this stage, and refining the model will have cascading effects on subsequent stages.

\subsection{Mechanoelectrical Transduction Model}

The next stage is the adaptive mechanoelectrical transduction (MET) model. As noted, in our study this is essentially a simplified and modified hair-cell transduction model \citep{Hudspeth2000-hx,Assad1992-fw,Shepherd1994-wh}. The components of this model are an adaptive motor with position $X_a$ and force $F_a$, a gating spring with stiffness $K_{gs}$, and an extent spring with stiffness $K_{es}$. Input is the scolopidia displacement $X_d$ and the output is an excitatory mechanotransduction current $I_{\mathrm{MET}}$. In a subsequent stage (see below), we transform this output into a firing rate to compare with our cell-attached recordings. 

We are not the first to use hair cell inspired models to attempt to recreate JO and JON responses \citep{Albert2007-km,Nadrowski2008-jt,Gopfert2006-ii,Gopfert2005-qd,Lee2018-dq,Lee2020-re,Effertz2012-fo}. This choice is sensible given the parallels to other animal auditory organs \citep{Gopfert2016-pl}, and given the relatively short response latencies which are observed from direct hair cell stimulation in vertebrates \citep{Corey1979-as}, and invertebrates \citep{Furukawa1967-jb} are also seen in \textit{Drosophila} following antennal displacement \citep{Albert2007-km,Lehnert2013-cl}. The goals of previous JON models were to explain AFP data and the antenna’s mechanical responses to sound stimuli or force steps. AFP responses are explained by the probabilistic opening of channels that follows the tension of the gating spring, while mechanical responses are explained by changes in compliance expected transduction channels opening, a so called “gating compliance” change. We discuss each of these goals below, and use them to highlight how our approach differs from previous models and how it might be extended.

To explain stimulus-driven changes in the AFP, previous models divide JONs into two groups, each responsive to linear, but sign opposed, motions of the antenna. They focus on the rapidly adapting JONs, the A and B types we targeted (but see \citep{Chang2016-lf} for an exception in a different model class). Like in our case, scolopidia are stretched by step stimuli, or the alternating phases of a vibrational stimulus, and tension on the gating spring activates mechanotransduction. These models do not rely on $I_{\mathrm{MET}}$ for a spiking model, but rather, they relate the AFP to the channel $P_o$. Following antennal displacement, adaptive motors re-equilibrate to new positions showing complete adaptation \citep{Nadrowski2008-jt}, following hair cell gating spring models \citep{Assad1992-fw}.

One challenge in previous studies is that they relied on a measure of the full population activity of JO. Any asymmetries present within the population may therefore need to be averaged over for simplification. As a likely consequence, there is no extent spring included to account for responses to static antennal displacement. Further, they make an assumption that resting $P_o = 0.5$, like some types of higher vertebrate hair cells \citep{Fettiplace2014-gu}. This is unlike the more general case for hair cells, where $P_o < 0.5$ \citep{Fettiplace2014-gu}, and which seems more consistent with observations in other insect auditory organs \citep{Hill1983-yu,Warren2018-ls}. Discrepancies aside, these models highlight the utility of using a motile element by showing adaptation to changes in antennal position, particularly with static offsets. 

While motor adaptation seems likely in \textit{Drosophila} JONs, its mechanism is unclear. The myosin motors responsible for adaptation in hair cells \citep{Gillespie1997-hd} have only been implicated in developmental roles in JONs \citep{Todi2008-hw}. Instead, the motors responsible for JON motility are thought to be dynein heavy-chain arms along the axoneme in the proximal region of the ciliated dendrite \citep{Gong2004-wu,Cheng2010-fc,Lee2010-bp}. However, the “9+2” axonemal structure normally implicated in ciliary motility (as in the flagella of sperm and \textit{Chlamydomonas} \citep{King2016-ew}) is absent. In its place is a “9+0” arrangement, meaning scolopidia axonemes lack a central pair of microtubules. Until recently, this observation has made motility somewhat controversial \citep{Yack2004-hr}. However, mutations effecting active motion generated by the antenna \citep{Gopfert2003-ao,Eberl2000-db}, recently implicated dynein heavy chains \citep{Karak2015-mx}, and the adaptation in the AFP strongly support motility.

From a modeling point of view, the dynamics of motility are clearly important, but we have made the simplifying assumption that this force is constant. In reality, $F_a$ may be $[Ca^{++}]$ dependent, as in hair cells \citep{Hacohen1989-cf}. In support of this idea, there are numerous centrin binding sites along axonemes \citep{King2016-ew} including those in insects \citep{Wolfrum1997-yj}. The precise arrangement of the transduction channel relative to the motor is unknown, although they are likely in the same proximal ciliary region \citep{Gong2004-wu,Cheng2010-fc,Zhang2013-gm}. Even direct $Ca^{++}$ feedback onto the motor however, does not require the motor to be located especially near the transduction channel. In the turtle hair cells, for instance, $Ca^{++}$ diffuses more than 100$nm$ to affect the adaptation motor \citep{Gillespie2009-va}.

Similarly, we have assumed the only variable governing $P_o$ is gating spring tension. As was the likely case for the adaptation motor, $Ca^{++}$ may provide feedback, in this case onto the channel itself or through a second-messenger system. The former possibility would be similar to the case in hair cells, where $Ca^{++}$ causes rapid adaptation \citep{Cheung2006-xm,Choe1998-kw,Ricci2000-dz}. In our simple model, $FR$ responses are often sluggish and do not terminate as rapidly as those suggested by the data. This could be due to mis-specification in the spiking model (discussed below); alternatively, $Ca^{++}$ feedback directly sculpting $I_{\mathrm{MET}}$ responses is an attractive mechanism which could sharpen responses considerably. 

In defining a potential role for $Ca^{++}$ it is important to note the composition of the scolopale lymph which bathes the cilia is not precisely known. It has been presumed to be similar to that of other insect mechanosensory organs \citep{Kuppers1979-qz,Thurm1974-xl}, where active transport creates a $K^+$ rich environment. Further support for this comes from phenotypes of $Na^+ / K^+ \mathrm{ATPase}$ knockdown in \textit{Drosophila} JO \citep{Roy2013-ii}, but $[Ca^{++}]$ is still unknown. One potentially useful technique to determine this value would use fluorescent ratiometric indicators. Loading dye indicators into JONs or scolopidia may be challenging, but recently developed ratiometric GECIs may provide an alternative route \citep{Cho2017-iq,Luo2019-ux}. A simpler first test for involvement would be to use sensitive, high SNR GECIs \citep{Zhang2020-yz} to see if $Ca^{++}$ invades the cilium during stimulation with 2-photon microscopy. 

Another major goal of previous models was to explain changes in stimulus-evoked antennal movement, and specifically the observation that the antenna’s compliance can increase. In these models, as channels “swing” open, their springs become displaced from the channel some distance $D_g$, and the stiffness decreases as the contribution of those springs to the overall organ stiffness lessens. Once enough channels are opened, however, many springs now share an offset distance from the channel, and the overall stiffness increases again. This transient softening (compliance change) has been observed in hair cells \citep{Howard1988-qf}, and is thought to be a general feature of bistable viscoelastic systems (e.g. muscle spindles).

Analogously to hair cell models, previous work on JON modeling has explicitly tied the compliance changes observed in the antenna to transduction channel gating compliance. These models successfully predict changes in stiffness resulting from rapid steps in applied force, which presumably results from coordinated gating within many JONs simultaneously \citep{Albert2007-km}. 

We have not incorporated gating compliance into our models. This was primarily a choice of convenience, as omitting $D_g$ significantly simplifies the model. Moreover, we observed that steady state, or simple motion was still reasonably well explained without this element. However, incorporating  hair-cell model $D_g$ dynamics could produce a better fit to our data, and in particular it could improve the fit to rapidly-fluctuating stimuli.

We did observe another type of adaptation in our experiments, particularly during vibrational stimuli. This adaptation appeared to leave the transduction channel closed, and could be recruited on a slower timescale than we expected for the transduction channel $Ca^{++}$ feedback. For this reason we believe this to be the result of a compliance change that lowered the tension on the gating spring. We proposed a stiff linker which becomes unbound in a force-dependent manner, resulting in a greater compliance. We think it is logical that the force-dependent compliance changes observed at the level of antennal mechanics is the same mechanism we observe during these responses. 

Many studies have reported responses which decrease as intensity increases \citep{Effertz2011-yt,Kamikouchi2009-mz,Tootoonian2012-qc}. During steps in the amplitude of a noise stimulus Clemens and colleagues \citep{Clemens2018-wp,Clemens2021-mx} have observed a similar adaptation to the change in variance, so-called “variance adaptation”. Although we do not have directly analogous stimuli, the time course of changes appears similar. 

How might this compliance change come about? There are many options. One attractive option implicates the TRPN1 channel NompC. Mutations of NompC greatly diminish JON currents recorded from the GFN \citep{Lehnert2013-cl}, but do not eliminate them, consistent with residual compliance after removing a stiff element. NompC is situated in an advantageous location to adjust compliance as well. It localizes to the distal cilia \citep{Lee2010-bp}, while Nan-Iav, the likely candidate for the majority of the mechanotransduction current \citep{Warren2018-ls,Zhang2013-gm} is located in the proximal cilia, near the presumed adaptive motors \citep{Gong2004-wu,Cheng2010-fc,Lee2010-bp}. NompC has prominent ankyrin repeats, as seen in the recently solved cryoEM structure \citep{Jin2017-js}. These repeats appear to convey forces through microtubule binding experimentally \citep{Zhang2015-ut,Liang2013-ai}, and in computational simulations \citep{Argudo2019-gz}. In fact, they have long been thought to be the gating spring proper \citep{Howard2004-it,Liang2013-ai}. We proposed a “stiff linker,” with unbinding and rebinding, which is likely a large oversimplification. If NompC were this element, then the dynamics of unbinding and rebinding may be complex, like the case for hair cell tip-links \citep{Mulhall2021-lj}, for example, which also display $Ca^{++}$ sensitivity.

Clearly more work is required to investigate the nature of this link. While NompC is a reasonable candidate, there are many other viscoelastic elements within the scolopidium. Tropomyosin and MAP-2, are both localized to the rigid scolopale rods, suggesting a potential role for $Ca^{++}$ in adjusting their stiffness \citep{Wolfrum1997-yj,Yack2004-hr}. Dendritic caps, in particular, are well positioned to adjust tension on the scolopale. Their material is often described as “spongy” in electron micrographs \citep{Field1998-hv} and they may also show significant force-dependent changes in compliance. (Prahlad et al., 2017) showed cap cells might tension slowly, and have multiple timescales of relaxation. Still, with the pre-existing work on molecular manipulations of NompC, it remains an attractive candidate for experimentation.For example, using the ankyrin repeat mutants generated by Zhang and colleagues \citep{Zhang2015-ut} to look for changes in adaptation time constants.

To further clarify many of these outstanding issues, an obvious technique to utilize is whole-cell voltage-clamp. Isolated transduction currents would be far easier to interpret than model fits to firing rate output. \citep{Warren2018-ls} successfully performed these experiments in larger chordotonal organ neurons, and we expect a small number of recordings to be possible in \textit{Drosophila}. In particular, now with a catalogue of stimulus-evoked spiking responses, assessing the relative amount of disruption, for example from dialysis during whole-cell recordings, is more straightforward. 

\subsection{Somatic Spiking Model}

The final stage in our model for JON responses is a simple spiking model. Like in some biophysical models of hair cells \citep{Neiman2011-fh}, we have considered $I_{\mathrm{MET}}$ separately from the dynamics at the soma. Although many interesting properties likely result from feedback between the two systems (e.g. \citep{Amro2014-um}), in our case, from the point of view of the soma, the cilia is simply a current delivery system.

As a first approximation, excitability in neurons can be described using “fast-slow” systems \citep{Izhikevich2010-tq}. To generate spiking behavior when driven with an external current, these systems have a fast excitatory process akin to voltage-gated sodium channels, and a slow recovery process which lumps potassium activation with sodium channel inactivation. This is in contrast to full conductance models in which all (or many) channel conductances are specified along with other cellular properties. In our experiments we did not have access to these conductances, so we used the fast-slow approximation. More specifically, we use  a quadratic integrate and fire model with a recovery variable, also known as an Izhikevich model \citep{Izhikevich2010-tq},. This model formalism only considers the dynamics of a small region of phase space which is “spike relevant,” and adds an after-hyperpolarization term. This makes incorporating experimental observations like resting potential easier, and in principle allows models to be fit in a straightforward manner. 

We chose this model because we noticed that it can easily generate an electrical model that acts as a weakly tuned band-pass or high-pass filter, similar to those seen in \textit{Calliphora} \citep{Gewecke1970-eb}. This suggests a spiking model with a passive dendritic compartment would be a good choice. We avoided using a threshold model like a leaky integrate and fire neuron, because we wanted greater control over filtering, and to see if adjusting the magnitude of an after-hyperpolarization parameter could recreate any of the responses seen. 

As there have been no prior single JON studies in \textit{Drosophila}, and we lacked membrane voltage data in most cases, this is an especially underconstrained stage of our model, and it can be greatly improved. Fortunately there are many straightforward ways to accomplish this. The easiest improvement that could be made is to fit Izhikevich models to data from current injection protocols during current-clamp recordings. This would allow us to better relate the modeled $I_{\mathrm{MET}}$ currents to $FR$ responses. Adding more recovery variables, analogous to incorporating additional conductances, might be required to achieve good performance. Although other intermediate models exist (for example, a two-variable conductance model $I_{\mathrm{Na,p}}+I_{\mathrm{K}}$, \citep{Izhikevich2010-tq}), a more satisfying extension might be towards using a full conductance model \citep{Shilnikov2012-ms}. 

With this goal in mind, an ideal set of experiments would involve pharmacologically isolating currents during voltage-clamp, potentially with mechanotransduction blocked \citep{Nesterov2015-io,Warren2018-ls}. It is known, for instance, that there are diverse $K^+$ channels expressed in photoreceptors \citep{Anderson1996-ir} another ciliated peripheral sensory neuron. Recent work using RNAseq and pharmacology has shown $K^+$ conductances sculpt processing later in the \textit{Drosophila} visual system \citep{Gur2020-zc}, and the lack of synaptic feedback onto JONs \citep{Kamikouchi2010-iz} makes bath application of pharmacological agents straightforward. Moreover, sequencing of the antenna, along with other \textit{Drosophila} organs, has recently been completed \citep{Li2021-vc}, potentially allowing a more targeted approach. 

Even in the absence of voltage-clamp experiments, there are many opportunities to improve models and understanding. For example, a distribution of values for $I_{\mathrm{MET}}$ would significantly ease modeling. The resonant properties of the neurons could also be investigated by current injection \citep{Hudspeth1988-rz,Crawford1981-op}, allowing, for example, time constants for $I_{K}$ to be better constrained. With pharmacology, contributions from different conductances can be further dissected. Of course, voltage-clamp “ZAP protocols” that sweep through frequencies of injected current to determine membrane resonance, would still be preferable \citep{Hutcheon1996-kd,Puil1986-fo}. 

Indeed, it would be interesting to see if electrical resonance played a role in specifying the frequency selectivity among different types of JONs, some of which have peak responses $>350Hz$. In the turtle cochlea this tuning requires BK-type conductances, while lower frequency tuning does not seem to depend as strongly on $Ca^{++}$ feedback, and rather relies on a large inward rectifying $K^+$ current \citep{Fettiplace1999-ap}. Tuning in JONs, and other \textit{Drosophila} chordotonal organ neurons may also use similar mechanisms, or perhaps more nontraditional ones \citep{Azevedo2017-kj}.

\section{Outlook}

We now briefly give an outlook on \textit{Drosophila} JON physiology. There have been many recent reviews in the fields of insect mechanotransduction \citep{Hehlert2020-xe,Eberl2016-yt}, mechanosensation \citep{Tuthill2016-ei}, and audition \citep{Kamikouchi2016-xf,Gopfert2016-pl,Hildebrandt2015-xc}. JONs are a key part of these fields, and a deeper understanding of their physiology will no doubt help guide our understanding of how auditory and mechanical signals are used by central circuits and guide adaptive behavioral responses.

Comprehensive physiology is, of course, elusive, but we are quite far from understanding the tuning across individual \textit{Drosophila} JONs. Our study has focused on A- and B-JONs, which themselves are known to have diverse tuning \citep{Ishikawa2017-ew,Patella2018-tk}. Beyond A- and B-JONs, what do other JON responses look like? Perhaps they are more similar to those observed in the \textit{Drosophila} femoral chordotonal organ \citep{Mamiya2018-iv}. Given the large ultrastructural diversity among chordotonal organ neurons \citep{Field1998-hv}, modeling other JON types might require adding additional mechanical components. Targeted recording from existing driver lines would be required to answer this question.

In terms of a molecular understanding of mechanotransduction, \textit{Drosophila} is well poised to continue making rapid progress. The recently completed “Fly Cell Atlas” contains a sequenced antenna \citep{Li2021-vc}, among other organs, and there exists a wealth of molecular tools for manipulating JONs \citep{Kernan2007-ug}. Advances using high-throughput and optimized \textit{in situ} hybridization \citep{Meissner2019-ra}, could help map the distributions of molecules within JO, including, with care, various channel isoforms. Moreover, recent progress in understanding how dyenins support motility might help clarify the nature of motor-based adaptation in “9+0” axonemes as well \citep{Lin2018-ln}.

What is general about chordotonal organ neurons and what is specific to \textit{Drosophila} JONs? We know that the relative size of \textit{Drosophila} JONs likely results in some differences from larger insect chordotonal organ neuron physiology. For example, we did not observe discrete dendritic spikes \citep{Hill1983-ro,Warren2018-ls}. Among lower vertebrates and mammals there is a large amount of diversity in biomechanical and electrical processing \citep{Fettiplace1999-ap}. For example, hair cells in turtle auditory papilla sometimes show spiking responses and have responses dominated by electrical tuning, while the mammalian cochlear inner hair cell has responses governed more by mechanical features of the organ. Among arthropods, or indeed within the various chordotonal organs within an animal, there may be a similar amount of diversity. In concert with molecular techniques, advances in whole-organ nanometer-scale imaging \citep{Kuan2020-tt}, we may soon be able to map ultrastructural diversity to molecular diversity among scolopidia. 

Computational modeling has also become more accessible in recent years. Advances in variational inference \citep{Salvatier2016-lh}, just-in-time compilers \citep{Lam2015-rf,Bezanson2017-lk}, and scalable computation has democratized access to model classes previously requiring specialist knowledge. Domain experts in biological sciences can now, with care, build Bayesian models constrained by the same types of variance observed in data \citep{Schroder2019-ps}. On the other end of computation, by using pre-existing datasets, machine learning approaches have been used to solve for ion channel dynamics \citep{Goncalves2020-my} and hypothesize structures \citep{Jumper2021-mv}.

In sum, as a result of this progress and, to some extent, the present study, we are well poised to vastly expand our understanding of Johnston’s organ neuron biomechanics, physiology, and molecular composition.
